\section{Encoding symbols in a cyclic finite coordinate}
\label{sec:encoding}

We now express the word constraints by finitely many translational tiling
equations on a common unknown set.  A symbol will be represented by the
dependence of a function on one of its input digits.  The present section
constructs tests forbidding specified simultaneous dependences; the next
section will force enough dependence to obtain an array. Encoding
functional constraints by tiling equations is developed in
\citet[Sections~4--8]{GTTwoTiles} and
\citet[Theorem~4.1]{GTPeriodic}. The issue here is to carry out the
required tests with a cyclic finite factor; the graph, dependence,
and cycle constructions below supply their exact local forms.

\subsection{Channels and the ambient group}

Retain the prime $p>200$, the width $N=p^2$, and the alphabet
$\Sigma=\Fp^\times$ from the preceding section.  Recall the line maps
$L_n(x)=x_2+nx_1$ and their kernel generators $u_n=(1,-n)$ from
\eqref{word:line-maps}.  The set of
\emph{channels} is
\[
 \mathcal I=\{0,\ldots,N-1\}\sqcup\{\eta_1,\eta_2\}.
\]
The first $N$ channels are ordinary; the last two are seeds.  Their input
digit labels are $J_n=\Sigma$ and $J_{\eta_t}=\{*\}$, respectively.
The seeds will eventually require the ordinary channels to display the
constant words $1$ and $2$ at some points.  Set
\[
 S=(\Zmod{p^2})^2,\qquad s(x)=x\bmod p^2,\qquad D=\abs{S}=p^4.
\]
Here $S$ records residues of the horizontal variable $x$.

For each $i\in\mathcal I$, choose primes $a_i,b_i$ and put
$P_i=\Zmod{a_i}\times\Zmod{b_i}$.  Require
\[
 a_{\eta_1}=2,\qquad a_{\eta_2}=3,\qquad
 b_{\eta_1},b_{\eta_2}>D,
\]
and choose all the $a_i,b_i$ to be distinct and different from $p$.
For each $j\in J_i$, choose a further prime $r_{ij}$, with all these
primes distinct and disjoint from the previous choices, such that
\begin{equation}
 r_{ij}=A_{ij}a_i+B_{ij}b_i,
 \qquad A_{ij},B_{ij}\in\Z_{\geq0}.
 \label{enc:digit-partition}
\end{equation}
These choices are possible using only finitely many primes.  Indeed, if
$a,b$ are coprime and $r\geq(a-1)b$, choose $B\in\{0,\ldots,a-1\}$
with $Bb\equiv r\pmod a$; then $A=(r-Bb)/a$ is a nonnegative integer.
There are arbitrarily large primes outside any prescribed finite set,
so this argument applies successively to every required $r_{ij}$.

Write
\[
 r_i=\prod_{j\in J_i}r_{ij},\qquad
 K_i=\Zmod{r_i}\cong\prod_{j\in J_i}\Zmod{r_{ij}},\qquad
 K=\prod_{i\in\mathcal I}K_i,\qquad
 P=\prod_{i\in\mathcal I}P_i.
\]
The displayed identification of $K_i$ is the Chinese remainder
isomorphism, given by taking residues.  Let $U_{ij}\leq K_i$ be the
coordinate subgroup supported on the digit $j$ in this identification.
We work in
\begin{equation}
 G=\Z^2\times V,\qquad
 V=\Zmod D\times\prod_{i\in\mathcal I}
       \bigl(\Zmod{r_i^2}\times P_i\bigr).
 \label{enc:ambient}
\end{equation}
The group $V$ is cyclic.  To see this, expand each $P_i$ into its two
cyclic factors.  The orders $D$, $r_i^2$, $a_i$, and $b_i$ of all the
resulting factors are pairwise coprime.  A tuple of generators therefore
has order equal to their product, which is $\abs V$.
The residue set $S$ is not an additional factor of $V$; the coordinate
of order $D$ in \eqref{enc:ambient} is the cyclic group $\Zmod D$.

\subsection{One graph and freely chosen high coordinates}

Write an element of $G$ as
\[
 \bigl(x,z,(v_i,c_i)_{i\in\mathcal I}\bigr),\qquad
 z\in\Zmod D,\quad v_i\in\Zmod{r_i^2},\quad c_i\in P_i.
\]
Its low coordinates are $k_i=v_i\bmod r_i\in K_i$.  Let
$q_0:G\to\Z^2\times K$ be the homomorphism retaining just $(x,k)$,
and let
\[
 H_0=\ker q_0
   =\{0\}\times\Zmod D\times
       \prod_{i\in\mathcal I}
       \bigl(r_i(\Zmod{r_i^2})\times P_i\bigr).
\]
Include the tiling equation
\begin{equation}
 A\oplus H_0=G.
 \label{enc:graph}
\end{equation}
For a target point in any fibre of $q_0$, every element of $A$ in that
fibre supplies exactly one representation with a summand in $H_0$.
Consequently \eqref{enc:graph} says precisely that $A$ contains one
point above each $(x,k)$.  We may thus denote its outputs by
$c_i(x,k)$, $z(x,k)$, and $v_i(x,k)$.

For subsequent constructions it is useful to describe $v_i$ by a
sheared high coordinate $\beta_i(x,k)\in K_i$, defined by
\begin{equation}
 v_i(x,k)=x_1+[k_i-x_1]_{r_i}+r_i\beta_i(x,k)
       \quad\text{in }\Zmod{r_i^2}.
 \label{enc:shear}
\end{equation}
Here $[\cdot]_{r_i}$ is the least nonnegative representative modulo
$r_i$.  For fixed $x,k_i$, the right side is a bijection from $K_i$
onto the elements of $\Zmod{r_i^2}$ with residue $k_i$.
Thus arbitrary functions $c_i$, $\beta_i$, and $z$ specify exactly one
graph satisfying \eqref{enc:graph}.  This is a choice of coordinates
on each fibre; it does not assert a direct-product splitting of
$\Zmod{r_i^2}$ into two groups of order $r_i$.

If we add $h=(h_1,h_2)$ to $x$ and $h_1$ to $v_i$, the new low
coordinate is $k_i+h_1$.  Since
\[
 [(k_i+h_1)-(x_1+h_1)]_{r_i}=[k_i-x_1]_{r_i},
\]
this operation preserves $\beta_i$.  The identity holds for every
integer lift of $h_1$, including negative ones.

\subsection{Dependence constraints on the useful outputs}

We require exactly the following restrictions on the functions $c_i$:
\begin{equation}
 \begin{aligned}
 c_n(x,k)&=\widetilde c_n(L_n(x),k_n)
       &&(0\leq n<N),\\
 c_{\eta_t}(x,k)&=\widetilde c_{\eta_t}(s(x),k_{\eta_t})
       &&(t=1,2).
 \end{aligned}
 \label{enc:dependence}
\end{equation}
The outputs $c_i$ will encode symbols; the other outputs will be used
to satisfy the activation equations.

\begin{lemma}\label{enc:invariance}
Fix a channel $i$ and a shift $\sigma\in\Z^2\times K$.
There is a finite nonempty tile $F_{i,\sigma}\subset G$ such that,
for every graph satisfying \eqref{enc:graph},
\[
 A\oplus F_{i,\sigma}=G
 \quad\Longleftrightarrow\quad
 c_i(b-\sigma)=c_i(b)\text{ for every }b\in\Z^2\times K.
\]
This equation imposes no restriction on the graph's other outputs.
\end{lemma}

\begin{proof}
Choose any lift $g\in G$ of $\sigma$ with $c_i$-coordinate zero, and
partition $H_0$ into
\[
 E_0^{(i)}=\{h\in H_0:c_i(h)=0\},\qquad
 E_{\ne0}^{(i)}=H_0\setminus E_0^{(i)}.
\]
Set
\[
 F_{i,\sigma}=E_{\ne0}^{(i)}\sqcup(g+E_0^{(i)}).
\]
The union is disjoint because its two parts have, respectively,
nonzero and zero $c_i$-coordinates.  Fix a target base $b$.
The unique graph point above $b$, together with $E_{\ne0}^{(i)}$,
covers exactly those points in that fibre whose $c_i$-coordinate
differs from $c_i(b)$, once each.  The graph point above $b-\sigma$,
together with $g+E_0^{(i)}$, covers exactly the points whose
$c_i$-coordinate is $c_i(b-\sigma)$, once each.
In both assertions, every choice of the other outputs and high
coordinates is attained exactly once, since those offsets range
freely in $E_0^{(i)}$.  The two contributions partition the target
fibre precisely when the displayed equality of useful outputs holds.
\end{proof}

Apply \cref{enc:invariance}, for each $i$, to a generator of every
$K_\ell$ with $\ell\ne i$, keeping all other base coordinates fixed.
This removes dependence on the other low inputs.  For an ordinary
channel $n$, also apply it to the horizontal shift $u_n$ with zero
low-coordinate shift.  Since $\ker L_n=\Z u_n$, the remaining
dependence on $x$ factors through $L_n(x)$.  For each seed use instead
the two horizontal shifts $(p^2,0)$ and $(0,p^2)$, again with zero
low-coordinate shift.  These give precisely dependence on $s(x)$.
There are only finitely many such generators and tiles, and their
equations are equivalent to \eqref{enc:dependence}.  In particular,
they leave all the functions $\beta_i$ and $z$ unrestricted.

For a graph with these dependences, write
$g_{i,x}:K_i\to P_i$ for its useful output at $x$.  Define its
\emph{active labels} by
\[
 \mathcal A_i(x)=
 \bigl\{j\in J_i:\text{there exist }w,w'\in K_i
   \text{ with }w-w'\in U_{ij}
   \text{ and }g_{i,x}(w)\ne g_{i,x}(w')\bigr\}.
\]
Thus a label is inactive exactly when changing that digit never
changes the output.  For ordinary $n$, the set $\mathcal A_n(x)$
depends only on $L_n(x)$; for a seed it depends only on $s(x)$.
A function independent of every digit is constant, since one can
pass between any two inputs by changing one digit at a time.
At this stage the active sets are allowed to be empty.

\subsection{A cycle test for simultaneous activity}

\begin{lemma}[Exclusion of simultaneous activity]\label{enc:exclusion}
Let $C=(i_1,\ldots,i_t)$ be a nonempty cyclically ordered list of
distinct channels, and choose a label $j(i)\in J_i$ for each $i\in C$.
There is a finite family of finite nonempty tiles such that a graph
satisfying \eqref{enc:dependence} tiles with every member of this
family if and only if, at every $x\in\Z^2$, at least one of the
chosen labels $j(i)$ is inactive.
\end{lemma}

\begin{proof}
For each $i\in C$, let $\operatorname{prev}(i)$ denote the preceding
channel in the cycle.  Choose arbitrary maps
\[
 d_i:P_{\operatorname{prev}(i)}\longrightarrow U_{i,j(i)}.
\]
For $e=(e_i)_{i\in\mathcal I}\in P$, put
\[
 d_i[e]=
 \begin{cases}
  d_i(e_{\operatorname{prev}(i)})&i\in C,\\
  0&i\notin C.
 \end{cases}
\]
Include a tile for every joint choice $d=(d_i)_{i\in C}$, namely
\begin{equation}
 F_d=\left\{
 \bigl(0,\zeta,(\upsilon_i,e_i)_{i\in\mathcal I}\bigr):
 \begin{array}{l}
 e\in P,\quad \zeta\in\Zmod D,\\
 \upsilon_i\in\Zmod{r_i^2},\quad
 \upsilon_i\bmod r_i=d_i[e]\quad(i\in\mathcal I)
 \end{array}
 \right\}.
 \label{enc:cycle-tile}
\end{equation}
This is an ordinary set: distinct $e$ give distinct useful-output
offsets, and every choice of $\zeta$ and the lifts $\upsilon_i$
occurs once.  Both the tile and the family of all maps $d$ are finite.

Fix a target base $(x,k)$.  For a fixed $e$, a representation using
\eqref{enc:cycle-tile} must start at the graph point over
$(x,(k_i-d_i[e])_i)$.  The useful outputs of the target are then
\begin{equation}
 \Phi_{x,k,d}(e)
   =\bigl(e_i+g_{i,x}(k_i-d_i[e])\bigr)_{i\in\mathcal I}.
 \label{enc:permutation}
\end{equation}
For any prescribed remaining coordinates of the target, the offsets
$\zeta$ and $\upsilon_i$ are uniquely determined: each low residue is
already correct, and all its lifts occur in the tile.
Hence $A\oplus F_d=G$ holds exactly when \eqref{enc:permutation} is a
permutation of $P$ for every $(x,k)$.  This criterion is independent
of the graph's high outputs and its $z$-output.

Suppose some $j(i_0)$ is inactive at $x$.  Then the $i_0$-component
of \eqref{enc:permutation} is $e_{i_0}+g_{i_0,x}(k_{i_0})$,
independent of the preceding component of $e$.
From any target useful-output tuple, we can therefore recover
$e_{i_0}$ uniquely.  The next equation around the cycle then recovers
the next component, and continuing recovers all components on $C$.
The starting equation remains satisfied because it never depended
on the preceding component.  Off the cycle, each equation is just a
fixed translation and recovers its component separately.  This gives
a unique preimage for every target, for all $k$ and all maps $d$.

Conversely, suppose every chosen label is active at some $x$.
For each $i\in C$ choose $w_i,w_i'\in K_i$ such that
\[
 w_i-w_i'\in U_{i,j(i)},\qquad
 \varepsilon_i:=g_{i,x}(w_i)-g_{i,x}(w_i')\ne0.
\]
Set $k_i=w_i$ on the cycle, and choose the other $k_i$ arbitrarily.
Extend $\varepsilon$ by zero off $C$ to an element of $P$.
For each $i\in C$, choose a map with
\[
 d_i(0)=0,\qquad
 d_i(\varepsilon_{\operatorname{prev}(i)})=w_i-w_i'.
\]
The two specified arguments are distinct, so these prescriptions
extend to maps on the whole domain.  They are among the maps tested
in the finite family.  At this base, the distinct inputs $0$ and
$\varepsilon$ have the same image under \eqref{enc:permutation}:
on $C$ their outputs are respectively $g_{i,x}(w_i)$ and
$\varepsilon_i+g_{i,x}(w_i')$, and off $C$ they coincide as well.
The permutation condition fails.  This proves both directions.
\end{proof}

Let $\mathcal W\subseteq\Sigma^N$ be the finite set of allowed words
from the word rule.  We impose the following specific instances of
\cref{enc:exclusion}:
\begin{enumerate}
 \item For every $w\in\Sigma^N\setminus\mathcal W$, use the cycle
 $(0,1,\ldots,N-1)$ with label $j(n)=w(n)$.
 \item For each $t\in\{1,2\}$, each ordinary channel $n$, and each
 $j\in\Sigma\setminus\{t\}$, use the two-channel cycle
 $(\eta_t,n)$ with labels $*$ and $j$.
\end{enumerate}
All these instances form a finite family.  Their precise consequences
are
\begin{equation}
 \prod_{n=0}^{N-1}\mathcal A_n(x)\subseteq\mathcal W,
 \qquad
 *\in\mathcal A_{\eta_t}(x)
   \ \Longrightarrow\ 
   \mathcal A_n(x)\subseteq\{t\}\quad(0\leq n<N).
 \label{enc:word-exclusions}
\end{equation}
The first product is identified with a set of words in the natural
order of the channels.  These statements also hold when some active
sets are empty.  The activation tiles constructed next will ensure
that every ordinary active set is nonempty at every point and that
each seed is active somewhere.

Every tile constructed in this section is a finite nonempty subset
of $G$.  As a check on the fibre counts, all have cardinality
\[
 \abs{H_0}=D\prod_{i\in\mathcal I}r_i\abs{P_i}:
\]
the dependence tiles partition a copy of $H_0$ into two shifted
pieces, while each cycle tile has $\abs P$ choices of $e$, $D$
choices of $\zeta$, and $r_i$ choices of each lift $\upsilon_i$.
